\documentclass{beamer}
\usetheme[numbering=progressbar]{focus}
\usepackage{tikz}
\usepackage{listings}
\usetikzlibrary{positioning}
\usetikzlibrary{shapes,arrows}
\usepackage{transparent}
\usepackage{fancyvrb}
\usepackage{listings}
\definecolor{main}{RGB}{47, 161, 219}
%\definecolor{textcolor}{RGB}{128, 128, 128}
\definecolor{background}{RGB}{240, 247, 255}
\definecolor{textcolor}{RGB}{85, 87, 83}
\title{Snake Oil Crypto:}
\subtitle{How I stopped to worry and started to love crypto}
\author{Jean-Louis Huynen}
\titlegraphic{\includegraphics[scale=0.20]{../../logos/d4-logo.pdf}}
\institute{Team CIRCL \\ }
\date{2019/12/06}

\begin{document}

\begin{frame}
    \maketitle
\end{frame}

\begin{frame}
        \frametitle{Outline}

        \begin{itemize}
          \item Use-Case: RSA,
          \item First Hands-on: Understanding RSA,
          \item Snake-Oil-Crypto: a primer,
          \item Second Hands-on: RSA in Snake-Oil-Crypto,
          \item D4 passiveSSL Collection,
          \item Interactions with MISP.
        \end{itemize}

\end{frame}

\begin{frame}
  \begin{center}
    {\bf Understanding RSA}
  \end{center}
\end{frame}

\begin{frame}
  \frametitle{RSA Basics}
  Ron {\bf R}ivest, Adi {\bf S}hamir, and Leonard {\bf A}dleman in 1977:
  \begin{itemize}
    \item asymmetric crypto system,
    \item can encrypt and sign,
    \item messages are big numbers,
    \item encryption is basically multiplication of big numbers,
    \item creates a \textit{trapdoor permutation}: turning x in y is easy, but
      finding x from y is hard.
  \end{itemize}

\end{frame}

\begin{frame}[fragile]
  \frametitle{RSA - Use with openssl}
  \begin{itemize}
  \item {\bf Hands-on}:

\begin{lstlisting}
~/hands-on/UsingRSA
\end{lstlisting}

   \item Decrypt message.bin
   \item generate a new private key,
   \item generate the corresponding public key,
   \item use this new key to encrypt a message,
   \item use this new key to decrypt a message. 
   
  \end{itemize}
\end{frame}


\begin{frame}[fragile]
  \frametitle{RSA ``by hand''}
run: sage rsa.sage at the folder's root:
\begin{lstlisting}[basicstyle=\tiny]
PlainText is: 1234567890
p = random_prime(2^32) = 2312340619
q = random_prime(2^32) = 2031410981
n = p*q = 4697314125248937239
phi = (p-1)*(q-1) = 4697314120905185640
e = random_prime(phi) = 2588085603940229747
d = xgcd(e,phi)[1] = -2102894211931680277
Does d*e == 1?
 mod(d*e, phi) = 1
CipherText y = power_mod(x, e, n) = 1454606910711062745
Decrypted CT is: 1234567890
\end{lstlisting}

\end{frame}

\begin{frame}
  \frametitle{With only one key}
  Several potential weaknesses:
  \begin{itemize}
    \item Key size too small: keys up to 1024 bits are breakable given the
      right means,
    \item close p and q,
    \item unsafe primes, smooth primes,
    \item broken primes (FactorDB, Debian OpenSSL bug).
    \item signing with RSA-CRT (instead of RSA-PSS)
  \end{itemize}

\end{frame}

\begin{frame}
  \frametitle{With a set of  keys}
  Several potential weaknesses:
  \begin{itemize}
   \item share moduli: if n1 = n2 then the keys share p and q,
   \item share p or q,
  \end{itemize}
 \vspace{10mm}
  {\bf In both case, it is trivial to recover the private keys.}
\end{frame}


\begin{frame}[fragile]
  \frametitle{Breaking small keys\footnote{https://www.sjoerdlangkemper.nl/2019/06/19/attacking-rsa/}}
  \begin{itemize}
\item {\bf Hands-on}:

\begin{lstlisting}
~/hands-on/SmallKey
\end{lstlisting}

   \item what is the key size of smallkey?
   \item what is n?
   \item what is the public exponent?
   \item what is n in base10?
   \item what are p and q?
   
  \end{itemize}

  \vspace{8mm}
  {\bf Let's generate the private key: }using p, then using q.
  
\end{frame}

\begin{frame}[fragile]
  \frametitle{Close Prime Factors}
  \begin{itemize}
\item {\bf Hands-on}:

\begin{lstlisting}
~/hands-on/ClosePQ
\end{lstlisting}

   \item use Fermat Algorithm\footnote{\url{http://facthacks.cr.yp.to/fermat.html}} to find {\bf both p and q:}

\begin{lstlisting}[basicstyle=\tiny]
def fermatfactor(N):
  if N <= 0: return [N]
  if is_even(N): return [2,N/2]
  a = ceil(sqrt(N))
  while not is_square(a^2-N):
    a = a + 1
  b = sqrt(a^2-N)
  return [a - b,a + b]
\end{lstlisting}

  \end{itemize}
 
\end{frame}

\begin{frame}[fragile]
  \frametitle{Shared prime factors}
     Researchers have shown that several devices generated their keypairs
   at boot time without enough entropy\footnote{Bernstein, Heninger, and Lange: \url{http://facthacks.cr.yp.to/}}:
   
\begin{lstlisting}[language=python, basicstyle=\tiny]
prng.seed(seed)
p = prng.generate_random_prime()
// prng.add_entropy()
q = prng.generate_random_prime()
n = p*q
\end{lstlisting}

Given n=pq and n' = pq' it is trivial to recover the shared p by computing their
{\bf Greatest Common Divisor (GCD)}, and therefore {\bf both private
  keys}\footnote{\url{http://www.loyalty.org/~schoen/rsa/}}.\\
\vspace{5mm}
``They cracked cracked about 13000 of them''
\end{frame}

\begin{frame}[fragile]
  \frametitle{Shared prime factors}
  \begin{itemize}
\item {\bf Hands-on}:

\begin{lstlisting}
~/hands-on/SharedPrimeFactor
\end{lstlisting}

\item Read README.txt, you have a challenge to solve :

  \begin{itemize}
  \item the \emph{answers} folder should be left alone for now,
  \item \emph{scripts} contains scripts that may be useful
    to solve the challenge,
  \item \emph{attempts} may hold your attempt are
    generating private keys. 
  \item \emph{bgcd-bd.sage} contains Daniel J. Berstein's algorithm for computing RSA
    collisions in batches.
  \end{itemize}

  \end{itemize}
 
\end{frame}


\input{soc-d4.tex}

\input{d4-pssl.tex}

\begin{frame}
  \frametitle{First release}
  \begin{itemize}
  \item[\checkmark] sensor-d4-tls-fingerprinting
    \footnote{\url{github.com/D4-project/sensor-d4-tls-fingerprinting}}:
    {\bf Extracts} and {\bf fingerprints} certificates, and {\bf computes} TLSH fuzzy hash.
  \item[\checkmark] analyzer-d4-passivessl
    \footnote{\url{github.com/D4-project/analyzer-d4-passivessl}}:
    {\bf Stores} Certificates / PK details in a PostgreSQL DB.
  \item snake-oil-crypto 
    \footnote{\url{github.com/D4-project/snake-oil-crypto}}:
    {\bf Performs} crypto checks, push results in MISP for notification
  \item lookup-d4-passivessl
    \footnote{\url{github.com/D4-project/lookup-d4-passivessl}}:
    {\bf Exposes} the DB through a public REST API.
  \end{itemize} 
\end{frame}

\begin{frame}
\frametitle{Get in touch if you want to join/support the project, host a passive ssl sensor or contribute}
\begin{itemize}
\item Collaboration can include research partnership, sharing of collected streams or improving the software.
\item Contact: info@circl.lu
\item \url{https://github.com/D4-Project} -  \url{https://twitter.com/d4_project}
\end{itemize}
\end{frame}

\end{document}
